% Copyright 2026  Ed Bueler

\documentclass[10pt,
               svgnames,
               hyperref={colorlinks,
                         citecolor=DeepPink4,
                         linkcolor=FireBrick,
                         urlcolor=blue},
               usepdftitle=false]{beamer}

\usetheme{Madrid}
\usecolortheme{beaver}
\setbeamercovered{transparent}
\setbeamerfont{frametitle}{size=\large}
\setbeamercolor*{block title}{bg=red!10}
\setbeamercolor*{block body}{bg=red!5}

\usepackage[english]{babel}
\usepackage[latin1]{inputenc}
\usepackage{times}
\usepackage[T1]{fontenc}
% Or whatever. Note that the encoding and the font should match. If T1
% does not look nice, try deleting the line with the fontenc.

\usepackage{fancyvrb,empheq,bm,xspace,dsfont,esint}
%\usepackage{minted}
\usepackage{tikz}
\usetikzlibrary{positioning,arrows.meta}

% If you wish to uncover everything in a step-wise fashion, uncomment
% the following command: 
%\beamerdefaultoverlayspecification{<+->}

\newcommand{\bb}{\mathbf{b}}
\newcommand{\bc}{\mathbf{c}}
\newcommand{\bbf}{\mathbf{f}}
\newcommand{\bl}{\bm{\ell}}
\newcommand{\bn}{\mathbf{n}}
\newcommand{\bq}{\mathbf{q}}
\newcommand{\br}{\mathbf{r}}
\newcommand{\bs}{\mathbf{s}}
\newcommand{\bx}{\mathbf{x}}
\newcommand{\by}{\mathbf{y}}
\newcommand{\bv}{\mathbf{v}}
\newcommand{\bu}{\mathbf{u}}
\newcommand{\bw}{\mathbf{w}}

\newcommand{\bE}{\mathbf{E}}
\newcommand{\bV}{\mathbf{V}}

\newcommand{\cB}{\mathcal{B}}
\newcommand{\cE}{\mathcal{E}}
\newcommand{\cM}{\mathcal{M}}
\newcommand{\cR}{\mathcal{R}}

\newcommand{\bzero}{\bm{0}}

\newcommand{\CC}{\mathbb{C}}
\newcommand{\NN}{\mathbb{N}}
\newcommand{\RR}{\mathbb{R}}
\newcommand{\ZZ}{\mathbb{Z}}

\newcommand{\one}{\mathds{1}}

\newcommand{\Div}{\nabla\cdot}
\newcommand{\grad}{\nabla}

\newcommand{\range}{\operatorname{range}}
\newcommand{\Span}{\operatorname{span}}

\newcommand{\Matlab}{\textsc{Matlab}\xspace}
\newcommand{\eps}{\epsilon}

\newcommand{\Lloc}{L_{\text{loc}}^1}

\newcommand{\ds}{\displaystyle}

\newcommand{\ip}[2]{\left<#1,#2\right>}

\newcommand{\twovect}[4]{\ensuremath{{#1}_{#2} =
                            \begin{bmatrix} #3 \\ #4 \end{bmatrix}}}

\newcommand{\rbullet}{{\color{FireBrick} \bullet}}

\newcommand*\circled[1]{\tikz[baseline=(char.base)]{
            \node[shape=circle,draw,inner sep=2pt] (char) {#1};}}

\newcommand{\textbook}{Saxe, \emph{Beginning Functional Analysis}}

\DefineVerbatimEnvironment{mVerb}{Verbatim}{numbersep=2mm,
frame=lines,framerule=0.1mm,framesep=2mm,xleftmargin=4mm,fontsize=\scriptsize}

\newcommand{\itemnum}[1]{\setcounter{enumi}{#1}\usebeamertemplate{enumerate item}}



\title[Poisson equation and Sobolev spaces]{the Poisson equation $- \grad^2 u = f$, \\ and the Sobolev spaces $H^k(\Omega)$}

\subtitle{calculations for weeks 12 \& 13}

\author{Ed Bueler}

\institute[]{UAF Math 617 Functional Analysis}

\date{Spring 2026}


\begin{document}
\beamertemplatenavigationsymbolsempty

\begin{frame}
  \maketitle
\end{frame}


\begin{frame}{Outline}
  \tableofcontents[hideallsubsections]
\end{frame}

\AtBeginSection[]
{% do not put toc here this time
}


\section{the Poisson equation: $- \grad^2 u = f$}

\begin{frame}{the Poisson equation}

\begin{itemize}
\item let $\Omega \subset \RR^d$ be open
\item you can assume $d=1,2,3$; I will usually draw pictures for $d=2$
\item assume the boundary of $\Omega$ is not too crazy (more below)
\item assume $f\in L^2(\Omega)$ is (square-)integrable

\begin{block}{Poisson equation}
   $$- \grad^2 u = f$$
\end{block}
\item picture:
\end{itemize}

\vspace{30mm}
\end{frame}


\begin{frame}{the Laplacian}

\begin{itemize}
\item given $f$, finding $u$ from the Poisson equation $- \grad^2 u = f$ is a problem of inverting the Laplacian operator

\begin{definition}[Laplacian]
for $\Omega\subset\RR^d$ open and $u:\Omega\to \RR$, the \emph{Laplacian} of $u$ is
\begin{equation*}
  \grad^2 u = \Div(\grad u) = \frac{\partial^2 u}{\partial x_1^2} +\dots + \frac{\partial^2 u}{\partial x_d^2}
\end{equation*}
\end{definition}

\item in $d=1,2,3$: \quad $\grad^2 u = u'', \quad \grad^2 u = u_{xx}+u_{yy}, \quad \grad^2 u = u_{xx}+u_{yy}+u_{zz}$
\end{itemize}

\begin{definition}[Laplace's equation $=$ potential equation]
\begin{equation*}
  \grad^2 u = 0,
\end{equation*}
and a solution of this equation is called \emph{harmonic}
\end{definition}
\end{frame}


\begin{frame}{tool: divergence theorem}

\begin{itemize}
\item $\grad^2$ usually appears in models because a flux is proportional to a gradient, and then a divergence will connect the interior behavior to a boundary integral of the flux
\item we will need the divergence theorem from the \href{https://bueler.github.io/fa/assets/slides/S26/week10.pdf}{weeks 10\&11 slides}:
   $$\int_\Omega \Div \bV\,dm = \int_{\partial\Omega} \bV\cdot \bn\,ds,$$
\item when using this, $\bV$ will be the gradient of a scalar function
\item the boundary integral will be a ``total flux of $\bV$ through the boundary''
\end{itemize}
\end{frame}


\begin{frame}{boundary-value problems for the Poisson equation}

\begin{itemize}
\item in the absence of any boundary conditions there is an infinite-dimensional space of solutions to $\grad^2 w = f$ on $\Omega$
\item if $f=0$ then this is the space of all \emph{harmonic functions} on $\Omega$:
    $$\grad^2 w= 0$$
\item the space of harmonic functions is $\infty$-dimensional
\item complex analysis is, essentially, the study of this space
\item but we want to have models which make \emph{one} prediction, so we will focus on \emph{boundary-value problems} that will pick-out one solution $u$
\end{itemize}

\begin{block}{Dirichlet problem for the Poisson equation}
for $f\in L^2(\Omega)$ and $g_D\in L^2(\partial\Omega)$, find $u:\Omega\to \RR$ so that
\begin{align*}
- \grad^2 u &= f \quad \text{ on } \Omega \\
u &= g_D \quad \text{ on } \partial \Omega
\end{align*}
\end{block}
\end{frame}


\begin{frame}{boundary-value problems for the Poisson equation}

\begin{itemize}
\item instead of $u$ itself, the gradient $\grad u$ could be set to known values along the boundary, which makes sense because the gradient is often proportional to a physical flux
\end{itemize}

\begin{block}{Neumann problem for the Poisson equation}
for $f\in L^2(\Omega)$ and $g_N\in L^2(\partial\Omega)$, find $u:\Omega\to \RR$ so that
\begin{align*}
- \grad^2 u &= f \quad \text{ on } \Omega \\
\grad u\cdot \bn &= g_N \quad \text{ on } \partial \Omega
\end{align*}
where $\bn$ is the outward unit normal vector field along $\partial\Omega$
\end{block}

\begin{itemize}
\item there are other ways to write the (scalar) directional derivative along the boundary:
	$$\grad u\cdot \bn = \frac{\partial u}{\partial n} = u_n$$
\end{itemize}
\end{frame}


\begin{frame}{applications}

\begin{itemize}
\item the Poisson equation can be used to model heat conduction, electrostatic potential, elastic membranes, equilibrium distributions for random walks, or many other physical phenomena
\end{itemize}

\begin{block}{application: temperature of a conductive solid}
\begin{itemize}
\small
\item Fourier's law for heat conduction in solids says that the heat flux is $\bq = -k \grad u$, where $u$ is the temperature, $\bq$ is the heat flux, and $k$ is the conductivity
\item a solid has a heat capacity $c$ (heat gained by given increase in temperature) and a mass density $\rho$ (mass per unit volume)
\item given a heat source $f$ within the domain, conservation of energy says
    $$c\rho\, \frac{\partial u}{\partial t} = - \Div\bq + f \quad \iff \quad \frac{d}{dt}\left(\int_\Omega c\rho\, u \,dx\right) = - \int_{\partial\Omega} \bq\cdot\bn\,ds + \int_\Omega f \,dx$$
\item let's assume steady state, so $d/d t=0$ and $\partial/\partial t=0$
\item if $k$ is constant then $- \Div\bq = - \Div(-k \grad u) = k \grad^2 u$
\item we derive Poisson's equation
    $$0 = k \grad^2 u + f$$
\item also called the (steady) \emph{heat equation} or \emph{diffusion equation}
\end{itemize}
\end{block}
\end{frame}


\begin{frame}{applications}

\begin{block}{application: deflection of an elastic membrane}
\begin{itemize}
\small
\item suppose $z=u(x,y)$ is a vertical displacement function over a domain $\Omega\subset \RR^2$
\item suppose that along FIXME
\item we derive Poisson's equation
    $$- \grad^2 u = f$$
\end{itemize}
\end{block}
\end{frame}


% at start of sections 2,3,... put outline
\AtBeginSection[] {
    \begin{frame}{Outline}
    \setbeamertemplate{section in toc}[sections numbered]
    \tableofcontents[currentsection]%,hideallsubsections]
    \end{frame}
}

\section{the Poisson equation in weak form}

\begin{frame}{bar}

FIXME multiply by $v$ and integrate by parts

FIXME Dirichlet and Neumann cases

\begin{itemize}
\item x
\end{itemize}
\end{frame}


\begin{frame}{bar}

FIXME localized $v$; hint at FE method

\begin{itemize}
\item x
\end{itemize}
\end{frame}


\section{the Sobolev spaces $H^k(\Omega)$}

\begin{frame}{bar}

FIXME follow Braess II.1 presentation (pages 28-31)

FIXME note $H^0(\Omega) = L^2(\Omega)$, and $H^k(\Omega) = W^{k,2}(\Omega)$ in some books; $W^{k,p}(\Omega)$ is the general case

\begin{itemize}
\item x
\end{itemize}
\end{frame}


\begin{frame}{bar}

FIXME $H_0^1(\Omega)$

\begin{itemize}
\item x
\end{itemize}
\end{frame}


\begin{frame}{bar}

FIXME Poincare-Friedrichs

\begin{itemize}
\item x
\end{itemize}
\end{frame}


\section{why does a Poisson solution exist?}

\begin{frame}{bar}

FIXME: minimization route

\begin{itemize}
\item x
\end{itemize}
\end{frame}


\begin{frame}{bar}

FIXME: Riesz representation theory route

\begin{itemize}
\item x
\end{itemize}
\end{frame}


\section{why is that solution smooth?}

\begin{frame}{bar}

FIXME: sketch how ellipticity means $u\in H^{2+k}(\Omega)$ if $f\in H^k(\Omega)$

\begin{itemize}
\item x
\end{itemize}
\end{frame}


\begin{frame}{Poisson on a disc}

FIXME Fourier series along boundary converts to smooth in the interior

FIXME relate to definition of $H^k([-\pi,\pi])$

\begin{itemize}
\item x
\end{itemize}
\end{frame}


\section{lecture content in weeks 12 \& 13}

\begin{frame}{lecture content in weeks 12 \& 13}

\begin{itemize}
\item most of this material is only from these slides

\begin{block}{to know from slides:}
\begin{itemize}
\item x
\end{itemize}
\end{block}

\begin{block}{to prove:}
\begin{itemize}
\item x
\end{itemize}
\end{block}

\item Assignment 7 is posted at \quad  \href{https://bueler.github.io/fa/index.html}{\texttt{bueler.github.io/fa}}
\end{itemize}
\end{frame}

\end{document}
